% --------------------------------------------------
\subsubsection{Диагональ зорилтот матриц}
% --------------------------------------------------

Диагональ зорилтот матриц нь хөрөнгө тус бүрийн өгөөжийн
түүврийн вариацыг хадгалж, харин хөрөнгүүдийн өгөөж хоорондын
ковариацыг тэг гэж үзнэ. Уг зорилтот матрицыг

\[
T_{\mathrm{diag}}
=
\operatorname{diag}(S)
=
\operatorname{diag}
(s_{11},s_{22},\ldots,s_{NN})
\]

гэж тодорхойлно. Матриц хэлбэрээр

\[
T_{\mathrm{diag}}
=
\begin{pmatrix}
	s_{11} & 0      & \cdots & 0\\
	0      & s_{22} & \cdots & 0\\
	\vdots & \vdots & \ddots & \vdots\\
	0      & 0      & \cdots & s_{NN}
\end{pmatrix}.
\]

Өөрөөр хэлбэл,

\[
(T_{\mathrm{diag}})_{ij}
=
\begin{cases}
	s_{ii}, & i=j,\\
	0,      & i\neq j.
\end{cases}
\]

Иймээс диагональ зорилтот матриц нь хөрөнгө тус бүрийн
өгөөжийн түүврийн вариацыг өөрчлөхгүй хадгалах боловч
хөрөнгүүдийн өгөөж хоорондын ковариацыг тэг болгон
хялбаршуулдаг.

Диагональ зорилтот матрицад суурилсан шугаман агшилтын
үнэлэгчийг

\[
\widehat{\Sigma}_{\mathrm{diag}}
=
(1-\widehat{\alpha}_{\mathrm{diag}})S
+
\widehat{\alpha}_{\mathrm{diag}}T_{\mathrm{diag}},
\qquad
0\leq\widehat{\alpha}_{\mathrm{diag}}\leq1,
\]

гэж тодорхойлно. Энд
$\widehat{\alpha}_{\mathrm{diag}}$ нь диагональ зорилтот
матрицад харгалзах агшилтын эрчим юм.


\paragraph{Агшилтын эрчмийн үнэлгээ.}

Диагональ зорилтот матриц нь түүврийн ковариацын матрицын
диагональ элементүүдийг өөрчлөхгүй хадгалдаг тул агшилт нь
зөвхөн диагоналийн бус элементүүдэд үйлчилнэ. Иймээс
түүврийн ковариацын матриц $S$ болон диагональ зорилтот
матриц $T_{\mathrm{diag}}$ хоорондын квадрат Фробениусын
зайг

\[
\widehat{d}_{\mathrm{diag}}^{\,2}
=
\left\|
S-T_{\mathrm{diag}}
\right\|_F^2
\]

гэж тодорхойлно. Уг зайг матрицын элементүүдээр
илэрхийлбэл

\[
\widehat{d}_{\mathrm{diag}}^{\,2}
=
\sum_{\substack{i,j=1\\i\neq j}}^{N}s_{ij}^{\,2}.
\]

Учир нь $S$ болон $T_{\mathrm{diag}}$ матрицын диагональ
элементүүд ижил тул тэдгээрийн ялгавар диагональ дээр тэг
байна.

Агшилтын эрчмийг үнэлэхийн тулд өмнөх хэсэгтэй адил

\[
z_{tij}
=
(y_{ti}-\overline{Y}_i)
(y_{tj}-\overline{Y}_j)
\]

гэж тэмдэглэе. Тэгвэл

\[
s_{ij}
=
\frac{1}{n}
\sum_{t=1}^{n}z_{tij}
\]

байна.

Түүврийн ковариацын матрицын элементүүдийн түүвэрлэлтийн
хэлбэлзлийг

\[
\widehat{\pi}
=
\frac{1}{n}
\sum_{i=1}^{N}
\sum_{j=1}^{N}
\sum_{t=1}^{n}
(z_{tij}-s_{ij})^2
\]

гэж тооцно.

Диагональ зорилтот матриц нь $S$ матрицын диагональ
элементүүдээс өөрөөс нь байгуулагддаг тул агшилтын эрчмийг
тооцоход диагональ элементүүдтэй холбоотой хэсгийг
$\widehat{\pi}$-ээс хасаж тооцно. Үүнийг

\[
\widehat{\rho}
=
\frac{1}{n}
\sum_{i=1}^{N}
\sum_{t=1}^{n}
(z_{tii}-s_{ii})^2
\]

гэж тодорхойлно. Энд

\[
z_{tii}
=
(y_{ti}-\overline{Y}_i)^2.
\]

Иймээс агшилтад хамаарах диагоналийн бус элементүүдийн
түүвэрлэлтийн хэлбэлзлийг

\[
\widehat{\pi}-\widehat{\rho}
\]

гэж авна.

Диагональ зорилтот матрицад харгалзах агшилтын эрчмийг

\[
\widehat{\alpha}_{\mathrm{diag}}
=
\max
\left\{
0,
\min
\left(
1,
\frac{\widehat{\pi}-\widehat{\rho}}
{n\,\widehat{d}_{\mathrm{diag}}^{\,2}}
\right)
\right\}
\]

гэж тооцно. Ингэснээр

\[
0
\leq
\widehat{\alpha}_{\mathrm{diag}}
\leq
1
\]

нөхцөл хангагдана.

Эцэст нь диагональ зорилтот матрицад суурилсан агшилтын
ковариацын матрицыг

\[
\widehat{\Sigma}_{\mathrm{diag}}
=
(1-\widehat{\alpha}_{\mathrm{diag}})S
+
\widehat{\alpha}_{\mathrm{diag}}T_{\mathrm{diag}}
\]

гэж тооцно.


\paragraph{Тооцоолох алгоритм.}

Диагональ зорилтот матрицад суурилсан агшилтын ковариацын матрицыг дараах үе шаттайгаар тооцно.

\begin{enumerate}
	\item Хөрөнгө тус бүрийн дундаж өгөөжийг тооцно.
	\[
	\overline{Y}_i
	=
	\frac{1}{n}
	\sum_{t=1}^{n}y_{ti},
	\qquad
	i=1,2,\ldots,N.
	\]
	
	\item Түүврийн ковариацын матрицыг байгуулна.
	\[
	s_{ij}
	=
	\frac{1}{n}
	\sum_{t=1}^{n}
	(y_{ti}-\overline{Y}_i)
	(y_{tj}-\overline{Y}_j),
	\]
	томьёогоор $S=(s_{ij})$ матрицыг тооцно.
	
	\item Диагональ зорилтот матрицыг байгуулна. Түүврийн ковариацын матрицын диагональ элементүүдийг хадгалж,
	\[
	T_{\mathrm{diag}}
	=
	\operatorname{diag}
	(s_{11},s_{22},\ldots,s_{NN})
	\]
	гэж байгуулна.
	
	\item Түүврийн ковариацын матриц болон зорилтот матрицын хоорондын квадрат Фробениусын зайг тооцно.
	\[
	\widehat{\gamma}_{\mathrm{diag}}
	=
	\left\|
	S-T_{\mathrm{diag}}
	\right\|_F^2
	=
	\sum_{\substack{i,j=1\\i\neq j}}^{N}s_{ij}^{\,2}.
	\]
	
	\item Түүврийн ковариацын матрицын элементүүдийн түүвэрлэлтийн хэлбэлзлийг тооцно. Эхлээд
	\[
	z_{tij}
	=
	(y_{ti}-\overline{Y}_i)
	(y_{tj}-\overline{Y}_j)
	\]
	гэж тооцоод,
	\[
	\widehat{\pi}
	=
	\frac{1}{n}
	\sum_{i=1}^{N}
	\sum_{j=1}^{N}
	\sum_{t=1}^{n}
	(z_{tij}-s_{ij})^2
	\]
	гэж олно.
	
	\item Диагональ элементүүдэд харгалзах хэсгийг тооцно.
	\[
	\widehat{\rho}
	=
	\frac{1}{n}
	\sum_{i=1}^{N}
	\sum_{t=1}^{n}
	(z_{tii}-s_{ii})^2,
	\]
	энд
	\[
	z_{tii}
	=
	(y_{ti}-\overline{Y}_i)^2.
	\]
	
	\item Агшилтын эрчмийг тооцно.
	\[
	\widehat{\alpha}_{\mathrm{diag}}
	=
	\max
	\left\{
	0,
	\min
	\left(
	1,
	\frac{\widehat{\pi}-\widehat{\rho}}
	{n\,\widehat{\gamma}_{\mathrm{diag}}}
	\right)
	\right\}.
	\]
	
	\item Агшилтын ковариацын матрицыг байгуулна.
	\[
	\widehat{\Sigma}_{\mathrm{diag}}
	=
	(1-\widehat{\alpha}_{\mathrm{diag}})S
	+
	\widehat{\alpha}_{\mathrm{diag}}T_{\mathrm{diag}}.
	\]
\end{enumerate}